\documentclass[12pt]{article}

\newcommand{\duedate}{-----}
\newcommand{\assignment}{Lecture 13B Summary: Interpretability in Deep Learning}
% Change the following to your name and UNI.
\newcommand{\name}{Aksel Kretsinger-Walters, adk2164}
\newcommand{\email}{}

% NOTE: Defining collaborators is optional; to not list collaborators, comment out the
% line below. Maximum of two collaborators per problem set.
%\newcommand{\collaborators}{Vig Vigerton (\texttt{UNI}), Alice Bob (\texttt{UNI})}
% No collaborators on PS 0

\makeatletter
\def\input@path{{../}{../../}{../../../}} % add as many parents as you need
\makeatother
\input{pset_template_NNDL.tex} %% DO NOT CHANGE THIS LINE

\begin{document}
\psetheader %% DO NOT CHANGE THIS LINE

\section{Overview}
Deep neural networks achieve impressive performance, yet their internal mechanisms remain difficult to interpret.
This lecture explores methods for understanding model behavior through visualization, gradients, and analysis of learned representations.

\section{Input Gradients}
The gradient $\partial L / \partial x$ reveals the model’s sensitivity to small input changes and can serve as a tool for interpretability or debugging.
However, raw gradients are often noisy and emphasize local perturbations rather than meaningful global structure.

\subsection{SmoothGrad}
To reduce noise, SmoothGrad (Smilkov et al.) averages gradients across noisy perturbations of the same image:
\[
		S = \frac{1}{N}\sum_{i=1}^{N}\frac{\partial L}{\partial x}(x + \epsilon_i), \quad \epsilon_i \sim \mathcal{N}(0, \sigma^2 I)
\]
This produces smoother, more globally interpretable saliency maps.

\section{Feature Visualization}
Visualization techniques attempt to understand what hidden units represent:
\begin{itemize}
		\item \textbf{Activation maximization:} find input images that strongly activate a given neuron or layer.
		\item \textbf{Layer-wise feature maps:} earlier layers capture edges and textures; deeper layers respond to higher-level concepts.
\end{itemize}
\textbf{Gradient ascent on images} can be used to iteratively modify inputs to amplify certain activations, as demonstrated in the \textit{Distill.pub} visualizations.

\section{Adversarial Examples}
Small, imperceptible perturbations can cause large changes in predictions.
Given an image $x$ labeled “cat,” we can compute:
\[
		x' = x + \epsilon \cdot \text{sign}\!\left(\frac{\partial L}{\partial x}\right)
\]
to make the model misclassify it (e.g., as “dog”).
This \textit{Fast Gradient Sign Method (FGSM)} reveals that networks rely on fragile, non-human features.

Adversarial examples are transferable across models and persist even in physical settings (e.g., printed images or modified road signs), posing serious security challenges.

\section{Deep Dream}
The \textit{Deep Dream} algorithm accentuates patterns already detected by a model by reinforcing activations during backpropagation:
\[
		x \leftarrow x + \eta \frac{\partial h_\ell}{\partial x}
\]
where $h_\ell$ denotes activations at layer $\ell$.
Repeated updates make features “grow” in the image, producing surreal, highly textured visualizations.

\section{Multimodal Interpretability}
CLIP (Contrastive Language–Image Pretraining) aligns text and vision embeddings, enabling text-guided exploration of visual concepts and revealing semantically meaningful activations across modalities.

\section{Sparse Representations}
\textbf{Superposition:} concepts are distributed across multiple neurons.
\textbf{Sparse coding hypothesis:} each concept is encoded by strong activation of a small subset of neurons. Two notions:
\begin{itemize}
		\item \textbf{Population sparsity:} only a few neurons fire for each concept.
		\item \textbf{Lifetime sparsity:} each neuron activates for few inputs.
\end{itemize}

\subsection{Sparse Autoencoders}
Adding an $\ell_1$ penalty or top-$k$ sparsity constraint promotes more interpretable latent features.
Recent work (e.g., \textit{Transformer Circuits}) applies sparse autoencoders to transformer activations, discovering “monosemantic” neurons tied to human-interpretable concepts.

\section{Circuits and Compositional Analysis}
Interpretability extends beyond individual neurons to multi-layer \textbf{circuits}—structured pathways representing invariances or hierarchical features.
For instance, neurons detecting “windows” and “wheels” combine to form a higher-level “car” detector, which then contributes to even more abstract semantic units.

\section{Key Takeaways}
\begin{itemize}
		\item Input gradients and visualization offer insight but require careful interpretation.
		\item Adversarial examples expose model brittleness and highlight the gap between human and machine perception.
		\item Sparse and compositional representations hint that deep networks organize knowledge hierarchically and efficiently.
\end{itemize}

\end{document}